# A positive-d terminal at a finite-ratio distance-four corner

25 September 2026. This is new round6 work. It uses the same coefficient-only
critical-center argument as `POSITIVE_D_FINITE_EDGE_TERMINAL.md`, with a
separate exact full-tail calculation. It does not cover ratio endpoints or
flat jet boundaries.

## 1. Hypotheses and conclusion

The physical nonnegative coefficient vector belongs to (B_c^9), where
\(B_c\) is the bounded valuation ring of a real closed extension of the
reals. The independent test belongs to the corresponding ring (B_s) of
another real closed extension. Suppose

\[
\lambda=0\quad\hbox{exactly},\qquad d_0>0,\quad\alpha_0=0.
\tag{1}
\]

All other coordinates may have arbitrary nonnegative residues or opening
rates. At this distance-four residue face write

\[
Q_{c_0}(U,V)=q_{20}U^2+q_{11}UV+q_{02}V^2
+q_{21}U^2V+q_{12}UV^2+q_{22}U^2V^2.
\]

For a real (t_0\ne0), assume

\[
q_{20},q_{02}>0,\quad
q_{20}t_0^4+q_{11}t_0^2+q_{02}=0,
\quad q_{21}t_0^4+q_{12}t_0^2>0.
\tag{2}
\]

Let the exact quadratic jet (J_4(t_0)) have kernel

\[
(w_0,x_0,y_0,1),
\]

and suppose its principal block in the variables ((W,X,Y)) is positive
definite. Equivalently the normalized kernel is simple and this specified
complement is positive. This real rank/complement condition is stated
explicitly; no unverified support census is used to remove it.

Then the original full slack is SOS in (B_c\otimes B_s) for every
independent test with

\[
\bar u=\bar v=0,\qquad
[\bar W:\bar X:\bar Y:\bar Z]=[w_0:0:0:1],
\tag{3}
\]

with the following un-divided normal quantities:

\[
X-vx_0,\quad Y-vy_0,\quad u-vt_0
\]

are the normal quantities handled below. **No asymptotic comparison or
extra smallness assumption on these three differences is needed.** They
automatically have zero residue in (3), and the positive Gram is proved
uniformly on all their independent rates. In particular this statement
does not assume that \(u/v\) is bounded or that its residue is \(t_0\): the
construction is an identity in the un-divided variables. The selected

\(t_0\) specifies a coefficient-side jet center, not a condition on an
independent mixed or pure-tested ratio.

## 2. Full polynomial and residue Hessian

The pure-tested amplitude (vY-Z) is a unit in (3). Normalize it to (-1),
so

\[
Z=1+vY.
\]

The whole opened alpha packet becomes (4\alpha(1+u^2)). Substituting

\[
X=vx,\quad Y=vy,\quad u=vt,\quad W=w,\quad V=v^2
\]

formally, on the coefficient side only, gives the exact full identity

\[
F=4\alpha+v^2K(x,y,t,w,V).
\tag{4}
\]

The complete (K) is polynomial. Its constant \(V\) coefficient is

\[
\begin{aligned}
K_0={}&4e_1x^2+4e_2(w-x-t)^2+4e_3(tw-y)^2+4e_4y^2\\
&+4e_5(tw-y-1)^2+4\beta(x+t)^2\\
&+2d\{(tw-y)^2+(x+t)^2-2tw\}+4\alpha t^2.
\end{aligned}
\tag{5}
\]

At the real residue the last term vanishes, and this is precisely
\(J_4(t)[w,x,y,1]\). The stipulated positive three-state complement makes
the \((w,x,y)\) Hessian positive. Its minimized scalar is

\[
\frac{16(q_{20}t^4+q_{11}t^2+q_{02})}{\det M_J(t)},
\tag{6}
\]

where \(M_J\) is that positive principal jet block. Physicality and (2)
make the numerator a positive multiple of \((t^2-t_0^2)^2\). Therefore
the full four-variable Hessian of \(K_0\) at

\[
q_0=(x_0,y_0,t_0,w_0)
\]

is positive definite. Its optimized next coefficient is

\[
\gamma=\partial_V K(q_0,0)
=\frac{16(q_{21}t_0^4+q_{12}t_0^2)}{\det M_J(t_0)}>0.
\tag{7}
\]

The determinant calculation is the full state Schur-complement calculation,
after scaling \(X,Y\) by \(v\). At the zero of (6), derivatives of the
positive state determinant do not affect the next coefficient. This
calculation proves the real sign in (7); no state-block inverse is inserted
into a tensor certificate.

## 3. Exact centers, including both signs

Solve \(
\nabla_qK(q_*(c,V),V)=0
\) by the Nash implicit function theorem, justified by the positive
four-Hessian. Let \(k(c,V)=K(q_*(c,V),V)\), and put

\[
\kappa=k(c,0).
\]

If \(\kappa\ge0\), take \(q_c=q_*(c,0),V_c=0\). If \(\kappa<0\), solve

\[
\nabla_qK(q_c,V_c)=0,\qquad
K(q_c,V_c)+V_c K_V(q_c,V_c)=0.
\tag{8}
\]

Its residue Jacobian has the positive four-Hessian and final pivot
\(2\gamma>0\), with zero lower-left block. Hence the nearby solution
exists in the coefficient field and satisfies \(V_c>0\), by the local
monotonicity of \(k(V)+Vk'(V)\) from its negative value at zero.

In the latter branch the pure coefficient scalar

\[
\mu=4\alpha+V_cK(q_c,V_c)
\tag{9}
\]

is nonnegative. It is the original physical slack at the actual
coefficient-side realization \(v=\sqrt{V_c}\),

\[
(u,X,Y,W)=(v t_c,v x_c,v y_c,w_c),\qquad Z=1+vY.
\]

Thus the sign argument is valid for arbitrary physical coefficients, not
only coefficient vectors already at contact.

## 4. Verified full-tail grouping without division by v

Define the un-divided normals

\[
N_1=X-vx_c,\quad N_2=Y-vy_c,\quad N_3=u-vt_c,
\quad N_4=v(W-w_c).
\tag{10}
\]

Write \(z=W-w_c\) while expanding the full polynomial. Its four exact
linear coefficients in \(N_1,N_2,N_3,z\) are

\[
vK_x,\quad vK_y,\quad vK_t,\quad v^2K_w.
\]

Accordingly the whole linear part is \(v\nabla_qK(q_c,V)\cdot N\).
Every monomial of the full remainder belongs to

\[
(N_1,N_2,N_3,vz)^2.
\tag{11}
\]

The direct checker allocates all 55 monomials of that complete remainder
into (11). Terms such as \(N_3^2z^2\) retain \(z^2\) as a bounded
infinitesimal coefficient. Terms with one \(z\) and one other normal have
the required factor \(v\), unless they already have two of the first
three normals. Terms with \(z^2\) and no other normals have \(v^2\).
Consequently no step substitutes \(z=N_4/v\).

The exact resulting normal Gram has residue

\[
\tfrac12\operatorname{Hess}_{(x,y,t,w)}K_0(q_0)\succ0.
\tag{12}
\]

Since \(\nabla K(q_c,V_c)=0\), the linear coefficient polynomial is
divisible by \(V-V_c\). Its contribution couples \(N\) to the scalar
generator \(V-V_c\) with an extra \(v\), hence a zero-residue cross
coefficient. The fixed-center univariate profile gives, by exact monic
division, the two identities

\[
\begin{array}{ll}
\kappa\ge0:&F=4\alpha+\kappa v^2+G^THG,
\quad G=(N_1,N_2,N_3,N_4,V),\\[2mm]
\kappa<0:&F=\mu+G^THG,
\quad G=(N_1,N_2,N_3,N_4,V-V_c).
\end{array}
\tag{13}
\]

The final Gram residue is the positive block (12) together with the
positive last entry \(\gamma\). The bounded positive-residue Gram lemma
proves it SOS, and the remaining pure nonnegative point scalars are
squares. Restoring the event normalization multiplies the roots by a
pure tested unit. This proves the claimed original-ring certificate for
every rate in (3).

## 5. Nonvacuity and audit status

At the physical residue

\[
c_0=(2,0,1,1,1,1,1,0,1)
\]

one has

\[
Q=144\{(U-V)^2+2UV(U+V)+U^2V^2\}.
\]

The corner center is \((x_0,y_0,t_0,w_0)=(-1/2,1/2,1,1)\), and
\(\gamma=6\). The full four-Hessian is

\[
\begin{pmatrix}
32&0&24&-8\\0&32&-24&-24\\24&-24&48&12\\-8&-24&12&32
\end{pmatrix}\succ0.
\]

`verify_positive_d_ratio_corner.py` reconstructs the full polynomial and
verifies **9 identities / 31 scalar entries**, the 61-term original
normalized parent, the 60-term profile, all 55 remainder monomials, and
the full residue Hessian. Its report is
`POSITIVE_D_RATIO_CORNER_CHECKS.json`.

The theorem is local in the residue state \(w_0\). Other real kernel
directions, a vanishing complement, ratio endpoints, and a vanishing
off-edge coefficient have not been included by approximation.
